\subsection{Related work}\label{sec:relwork}

\paragraph{The calculus of relations: history and expressiveness.}
The calculus of relations (\CoR) predates first-order logic as a formal
system. De Morgan introduced relational composition and converse in his
work on the syllogism in the 1850s and 1860s
\citep{demorgan1856,demorgan1864}; Peirce developed the logic of
relatives~\citep{peirce1870,peirce1880}, and Schr\"oder devoted a volume
to it in 1895~\citep{schroeder1895}. This was before first-order logic
was isolated in Hilbert's lectures and Hilbert and Ackermann's
textbook~\citep{hilbert-lectures,hilbert-ackermann1928,moore1988}.
The central expressiveness problem was whether every quantified statement
about relations could be \emph{condensed} into this variable-free calculus.
Schr\"oder claimed that it could~\citep[p.~551]{schroeder1895}, but
Korselt found a counterexample, published by L\"owenheim in
1915~\citep{loewenheim1915,loewenheim1967}. In modern terms, the \FOfour\ sentence
$
\forall a_1\,a_2\,a_3\,a_4.\bigvee_{i,j; i\neq j} a_i = a_j
$
states that the domain has at most three elements, but is expressible
neither in \FOthree\ nor by a \CoR\ equation.

Tarski raised the expressiveness question in 1941~\citep{tarski1941}
and, according to Tarski and Givant, soon obtained the three-variable
characterization: \CoR\ defines exactly the \FOthree-definable binary
queries. A full proof appeared only in their 1987
monograph~\citep{tarski-givant1987}; see
also~\citep{maddux2006}. For historical accounts, see
\citep{pratt1992,maddux1991,burris-legris,brady2000}.
Their original recursive translation~\citep[\S3.9]{tarski-givant1987}
did not analyze its size. A straightforward bottom-up translation using
DNF conversion~\citep{hilbert-ackermann1928,kleene1952,schoening1989}
and quantifier elimination by relational composition gives a
single-exponential upper bound (Appendix~\ref{app:classical-translation});
the reverse \CoR-to-\FOthree\ translation is linear.

\paragraph{Relational computation and quantifier elimination.}
The operations underlying \CoR\ have long served as computational
primitives, often in calculi distinct from full \CoR.
Copilowish gave a matrix-based treatment of the calculus over finite
domains~\citep{copilowish1948}, while Levien and Maron applied relational
operations to retrieval and inference
\citep{levien-maron1965,levien-maron1967,levien1969}.
Codd's relational algebra laid the foundations of database querying
\citep{Codd1970,Codd1972}, followed by work on relational-expression
equivalence, optimization, and efficient join evaluation
\citep{aho-sagiv-ullman1979equiv,aho-sagiv-ullman1979opt,
SelingerEtAl1979,yannakakis1981}.
Binary relational-expression evaluation became an algorithmic subject
\citep{hunt-szymanski-ullman1974,hunt-szymanski-ullman1977,
szymanski-ullman1975,szymanski-ullman1977}, with applications to parsing
and recursive database queries
\citep{sippu-soisalon-soininen1985,grahne-sippu-soisalon-soininen1987,
grahne-sippu-soisalon-soininen1991}.
Related algebras also appear in relational programming, temporal
reasoning, database execution, and graph querying
\citep{maclennan1981,montanari1974,allen1983,boncz-kersten1999,
cruz-mendelzon-wood1987,FletcherEtAl2015,HellingsEtAl2021,
HellingsEtAl2022}.

A second line of work exploits relational composition as existential
quantifier elimination, using matrix multiplication for fast evaluation.
Early examples include transitive closure
\citep{furman1970,fischer-meyer1971,munro1971,oneil-oneil1973,
galil-margalit1993} and algorithms for triangles, cliques, cycles,
and small subgraphs
\citep{itai-rodeh1978,nesetril-poljak1985,alon-yuster-zwick1997,
kloks-kratsch-mueller2000,eisenbrand-grandoni2004,yuster-zwick2004}.
Already in 1977, Chandra and Merlin studied generalized join--project
operations and their evaluation by rectangular Boolean matrix
multiplication~\citep{ChandraMerlin1977}.
Williams's algorithm for prenex first-order graph properties
\citep{Williams2014} can be viewed in the same algebraic framework:
DNF-based quantifier elimination yields ordinary \CoR\ composition
for three quantifiers and typed composition over tuples for larger
prenex formulas, evaluated using rectangular matrix multiplication
(as in \citep{ChandraMerlin1977}).
Gao and Impagliazzo\ \citep{GaoImpagliazzo2019} also treated general
\FOthree\ formulas, beyond prenex normal form, using matrix-based
evaluation.
Related techniques appear in first-order graph algorithms
\citep{BringmannFischerKunnemann2019} and database join--project
and join--aggregate evaluation
\citep{DeepHuKoutris2020,hu-yi2020,Hu2024,AboKhamisHuSuciu2025}.
Most directly, our earlier work used the \FOthree-to-\CoR\ translation
for truly subcubic \FOthree\ model checking
\citep{nakamura-uezato2027}.

\paragraph{Succinctness results.}
Nakamura proved exponential lower bounds for positive \CoR\ and
$\Sigma_1\CoR$, even for circuits, and a superpolynomial separation over
one binary relation~\citep[\S6]{nakamura2022}. The case of full \CoR\
remained open (see \S\ref{sec:natural-approaches}). In our earlier
work~\citep[Corollary~40]{nakamura-uezato2027}, we ruled out
polynomial-time \FOthree-to-\CoR\ translations under the 3-OV
hypothesis, but not the existence of small equivalent representations.
Theorem~\ref{thm:main} gives an unconditional exponential lower bound
even for \CoR\ circuits, while Theorem~\ref{thm:fixed-vocabulary}
gives tight subexponential circuit bounds over one binary relation.

On finite linear orders with no additional relations, Grohe and Schweikardt proved a polynomial-size translation from \FOthree\ to \FOtwo\ sentences~\citep[Theorem~4.3]{grohe-schweikardt2005}. Since \FOtwo\ translates linearly into \CoR, this also yields polynomial-size \CoR\ equations. By contrast, our exponential succinctness gaps from \FOthree\ to both \FOtwo\ and \CoR\ hold even for sentences over general finite structures.
%
They also showed that full first-order logic is exponentially more
succinct than \FOthree\ on the same orders, already witnessed by
\FOfour~\citep[Theorem~5.3]{grohe-schweikardt2005}. The short \FOfour\
formulas combine divide-and-conquer for distances~\citep{savitch1970}
with the classical quantifier-based abbreviation trick, which avoids
duplicating recursive subformulas
\citep{fischer-rabin1974,stockmeyer1974,ferrante-rackoff1979,compton-henson1990,vorobyov1996,stockmeyer-meyer2002}.
The \FOthree\ lower bound draws on Adler--Immerman formula-size
games~\citep{adler-immerman2003}. However, the resulting distance
predicates admit linear-size \CoR\ circuits by repeated squaring,
$S_{i+1}=S_i;S_i$; our lower bound holds even with such sharing.

Vinall-Smeeth established exponential succinctness gaps from
$\mathrm{FO}^{\ell+1}$ to $\mathrm{FO}^{\ell}$ for every $\ell\ge 3$,
using higher-arity (non-binary) relations~\citep[Theorem~2]{vinall-smeeth2024}.
Our result establishes the \FOthree-to-\FOtwo\ case over binary
signatures, with an exponential lower bound on the number of
quantifiers. The linear \FOtwo-to-\CoR\ translation was already
known~\citep{nakamura2022}, while Proposition~\ref{prop:reverse}
gives an exponential gap in the reverse direction.

\paragraph{Lower-bound methods.}
As explained in \S\ref{sec:natural-approaches}, the witness-crossing
argument underlying Nakamura's formula-size-game proof for $\Sigma_1\CoR$
\citep{adler-immerman2003,nakamura2022} relies on embeddings, which fail
under unrestricted complement. We also considered other standard
lower-bound techniques, but their usual resource measures do not directly
fit our setting. Formula-size games
\citep{grohe-schweikardt2005,french2013,hella-vilander2019},
Ehrenfeucht--Fra\"{\i}ss\'e and pebble games
\citep{fraisse1954classif,ehrenfeucht1961,immerman1982,libkin2004},
and multi-structural games
\citep{immerman1981,fagin2021,carmosino2024} address formula size,
quantifier depth, variable count, or quantifier number, rather than the
number of distinct composition gates in circuits with sharing. Indeed,
$G_k$ has only one quantifier and admits a \CoR\ expression of
composition depth one; the issue is the \emph{number} of compositions.
Monotone-circuit lower bounds~\citep{razborov1985,alon-boppana1987}
are natural candidates for our positive query, but equivalent full-\CoR\
circuits may use unrestricted complement. Switching-lemma and
random-restriction methods~\citep{hastad1986} typically exploit
shallow Boolean-circuit structure, whereas we impose no restrictions on
Boolean nesting and charge only compositions.

Communication complexity suggests further approaches. Rank,
discrepancy, and Karchmer--Wigderson techniques
\citep{mehlhorn-schmidt1982,kushilevitz-nisan1997,karchmer-wigderson1990,jukna2012}
do not directly bound compositions when Boolean operations are free.
For example, set disjointness has rank $2^k$, yet its $1$-entries form
the complement of a union of only $k$ rectangles. Alternating existential
and universal witness tests also resemble the nondeterministic and
co-nondeterministic levels of the communication polynomial hierarchy
\citep{babai-frankl-simon1986,goos-pitassi-watson2018}, but we know of
no direct transfer of lower bounds to our circuit model. These approaches did not yield the required composition-count bound.

Our proof instead draws on the affine-plane and random colorings of
Jipsen, Maddux, and Tuza~\citep{jipsen-maddux-tuza1995}, which, as
explained in the technical overview, already yield the two-witness
property needed here. The new ingredient is our fold-based argument:
it bounds how many local folds each distinct operation can notice,
even with unrestricted complement and circuit sharing.
